\chapter{Polarisation in McStas}
\label{c:polarization}
\index{Polarisation|textbf}
\index{Neutron polarisation|see{Polarisation}}

% ---------------------------------------------------------------------
% Notation used throughout this chapter. These macros are local to this
% file (rather than added to the shared preamble) since they are
% specific to the polarisation formalism below and not used elsewhere
% in the manual. Bold/hatted vector and operator notation follows
% Lovesey (\cite{lovesey84}) and Williams (\cite{gavin}).
\providecommand{\PB}{\mathbf{P}}          % polarization vector
\providecommand{\tP}{\hat{\mathbf{P}}}    % unit polarization vector
\providecommand{\SB}{\mathbf{S}}          % simulated (McStas) polarization vector
\providecommand{\sB}{\mathbf{s}}          % neutron spin vector
\providecommand{\BB}{\mathbf{B}}          % magnetic field vector
\providecommand{\nB}{\mathbf{n}}          % arbitrary unit direction vector
\providecommand{\muB}{\boldsymbol{\mu}}   % magnetic moment vector
\providecommand{\muno}{\hat{\boldsymbol{\mu}}} % neutron magnetic moment operator
\providecommand{\tauB}{\boldsymbol{\tau}} % torque vector
\providecommand{\dB}{\mathbf{d}}          % intra-cell atom position vector
\providecommand{\lB}{\mathbf{l}}          % unit-cell lattice vector
\providecommand{\RB}{\mathbf{R}}          % atomic position vector (R_ld = l+d)
\providecommand{\Io}{\hat{\mathbf{I}}}    % nuclear spin operator
\providecommand{\so}{\hat{\mathbf{s}}}    % neutron spin operator
\providecommand{\sigmao}{\boldsymbol{\hat\sigma}} % vector of Pauli operators
\providecommand{\sigmaH}{\hat\sigma}      % single Pauli operator component
\providecommand{\rhoo}{\hat\rho}          % density matrix operator
\providecommand{\alphao}{\boldsymbol{\alpha}} % spin-dependent potential amplitude
\providecommand{\betao}{\boldsymbol{\beta}}   % spin-independent potential amplitude
\providecommand{\Q}{\mathbf{Q}}           % scattering vector
\providecommand{\tQ}{\hat{\mathbf{Q}}}    % unit scattering vector
\providecommand{\tN}{\hat{\mathbf{N}}}    % unit vector along atomic (ferromagnet) spin
\providecommand{\FN}{F_N}                 % unit cell nuclear structure factor
\providecommand{\FM}{F_M}                 % unit cell magnetic structure factor
\providecommand{\Ru}{R_\uparrow}          % spin-up reflectivity
\providecommand{\Rd}{R_\downarrow}        % spin-down reflectivity
\providecommand{\nup}{n^\uparrow}         % number of spin-up neutrons
\providecommand{\nd}{n^\downarrow}        % number of spin-down neutrons
\providecommand{\Pu}{P^\uparrow}          % probability of spin-up
\providecommand{\Pd}{P^\downarrow}        % probability of spin-down
\providecommand{\chiU}{\chi_\uparrow}     % spin-up eigenfunction
\providecommand{\chiD}{\chi_\downarrow}   % spin-down eigenfunction
\providecommand{\madsq}{\overline{|F_N(\Q)|^2}}       % mean-square structure factor (coherent)
\providecommand{\sqmad}{\left|\overline{F_N(\Q)}\right|^2} % square of the mean (site-incoherent complement)
\providecommand{\bd}{\overline{|B_{ld}|^2}}            % mean-square isotope/spin-incoherent term
% ---------------------------------------------------------------------

\MCS provides a set of components for simulating neutron beam polarisation
and the optical elements that manipulate it: polarising monochromators,
mirrors and guides, magnetic-field regions (including spin flippers and
spin-echo coils), and polarisation-sensitive monitors. This chapter
originated as an appendix written by P. Christiansen documenting the first
polarisation components (2006); it is now promoted to a full chapter and
brought up to date with the current \MCS~3.x component set and the
\texttt{pol-lib} field/precession framework that underlies it.

We first describe the polarisation vector and how it is related to the
neutron wave-function (section~\ref{sec:pol}) and then the physics of the
scattering processes relevant to polarisation (section~\ref{sec:scat}).
Section~\ref{sec:precession} describes in detail how \MCS propagates the
polarisation vector through magnetic fields -- the numerical precession
algorithm and the field functions available -- since this underlies every
magnetic-field component in the library. Section~\ref{sec:new} then
documents the current polarisation components themselves, organised by
role, and section~\ref{sec:test} lists example and test instruments,
including a worked spin-echo example.

We rely heavily on the books~\cite{lovesey84,gavin} for the underlying
physics, where the detailed calculations can be found.

The notation used here (and in~\cite{gavin}) is $P$ (scalar), $\PB$
(vector), $\tP$ (unit-vector), $\sigmaH$ (operator), and $\sigmao$
(vector of operators).

\section{The Polarization Vector}
\label{sec:pol}

The spin of the neutron is represented by an operator $\so$ for which only a
single component can be measured at one time. Each single measurement will
give a value $\pm 1/2$, but if we could make a large number of measurements on
the same neutron state, in each of the three axis directions, and then make
the average we get $\langle \so \rangle$. The polarization vector, $\PB$, is
then defined as:
\begin{equation}
  \label{eq:pol_single}
  \PB = \frac{\langle \so \rangle}{s},
\end{equation}
so that $-1 \leq |\PB| \leq +1$

For a neutron beam which contains $N$ neutrons, each with a
polarization $\PB_i$, the beam polarization is defined as:
\begin{equation}
  \label{eq:pol_beam}
  \PB = \frac{\sum_i \PB_i}{N}.
\end{equation}

If we have one common quantization direction (e.g. a magnetic field
direction) each neutron will either be spin up, $\uparrow$, or spin down,
$\downarrow$, and the polarization can be expressed as:

\begin{equation}
  \label{eq:pol}
  P = \frac{\nup-\nd}{\nup+\nd},
\end{equation}

where $\nu$ ($\nd$) is the number of neutrons with spin up
(down).

For a given neutron the probability of the neutron being spin up, $\Pu$, is:

\begin{equation}
  \label{eq:prob_spinup}
  \Pu = \frac{\nup}{\nup+\nd} = \frac{\nup + (\nd-\nd)/2}{\nup+\nd}
  = \frac{1+P}{2},
\end{equation}

and $\Pd = 1-\Pu = (1-P)/2$.

The expectation value of the 'spin' operator, $\sigmao$, which can be
expressed by the Pauli matrices, is the polarization vector $\PB$, $\PB =
\langle \sigmao \rangle \equiv \langle \chi | \sigmao | \chi \rangle$.  The
most general form of the spin wave-function $\chi$ for a neutron (spin 1/2)
is:

\begin{equation}
  \label{eq:neutron_wave}
  \chi = a\chi_\uparrow + b\chi_\downarrow,
\end{equation}

where $\chi_\uparrow$ and $\chi_\downarrow$ are eigenfunction of
$\hat{\sigma}^z$, and the complex coefficients $a$ and $b$ satisfy
$|a|^2 + |b|^2 = 1$.

By calculation we find:
\begin{eqnarray}
P_x & = & \langle \chi | \sigmaH_x | \chi \rangle
= 2 \text{Re}(a^\ast b) \\
P_y & = & \langle \chi | \sigmaH_y | \chi \rangle
= 2 \text{Im}(a^\ast b) \\
P_z & = & \langle \chi | \sigmaH_z | \chi \rangle
= |a|^2-|b|^2
\end{eqnarray}

This shows the relation of the polarization vector to the neutron wave
function.

The neutron magnetic moment operator can be expressed in terms of $\sigmao$,
as:
\begin{equation}
  \label{eq:magnetic}
  \muno = \mu_n \sigmao,
\end{equation}
which, as shown above, is related to the polarization vector.  \\

In our simulation we represent the polarization by the vector $\SB = (s_x,
s_y, s_z)$ which is propagated through the different components so it has the
correct relative orientation in each component. The probability for the spin
to be parallel a given direction $\mathbf{n}$ is then:

 \begin{equation}
   \label{eq:probmcstas}
   P(\uparrow|\nB) = \frac{1+\nB \cdot \SB}{2}.
 \end{equation}

This equation (from~\cite{pol_seeger}) is easy to understand. The
average spin along $\nB$ is $\nB \cdot \SB$ and the probability then
follows from Eq.~\ref{eq:prob_spinup}.

For an unpolarized beam, $\mathbf{S} = \mathbf{0}$ and all directions
are equally probable (50~\%).

Note that in our approach we do not decide if the neutron is up or down after
a given component, but instead keep track of as much information for as long
as possible.

In the following we will use $\PB$ to denote the polarization
vector. The most important variables used are:

\begin{tabular}{ll}

  $\Q$  & Scattering vector. \\
  $\PB$  & Polarization before a component (ingoing). \\
  $\PB_\perp$  & Polarization perpendicular to scattering vector,
  $\PB_\perp = \tQ \times (\PB \times \tQ$). \\
  $\PB'$ & Polarization after a component (outgoing). \\
  $\tN$ & Unit vector in direction of atomic spin ($\tN \cdot \tilde{\BB} = -1$ for a ferromagnet). \\
  $\FN$ & Unit cell nuclear structure factor. \\
  $\FM$ & Unit cell magnetic structure factor. \\
\end{tabular}
\\

The unit cell nuclear structure factor is defined as:

\begin{equation}
  \FN = \sum_\dB
  \exp(i\Q \cdot \dB)\overline{b}_d,
\end{equation}

where the $\dB$ is the position of the d'th atom within the unit cell,
and $\overline{b}_d$ is the average of $b_d$. In the simple case of a
single atom Bravais crystal one finds $\FN = \overline{b}$.

The unit cell magnetic structure factor is useful when the atoms in
the crystal only have spin orbital angular momentum, and simple when
the magnet is saturated (all spins are parallel or anti-parallel to
\emph{one} direction, $\sigma_d=\pm1$). It is then given as:

\begin{equation}
  \FM =
  \gamma_n r_0 \sum_\dB \exp(i\Q \cdot \dB)\frac{1}{2}g_d F_d(\Q)\langle
  \hat{S}_d\rangle \sigma_d,
\end{equation}

where $r_0=\frac{\mu_0}{4\pi}\frac{e^2}{m_e} = 2.818 \times 10^{-15}$m, $g=2$ is the
Land{\'e} splitting factor, and $F_d(\Q)$ is the magnetic form factor, which
is the Fourier transform of the magnetization density (normalized so that
$F_d(0) = 1$), and $\langle \hat{S}_d\rangle$ is the thermal average of the
ordered atomic spin.

In the following the Debye-Weller factor ($\exp(-W_d)$) have been ignored in
all cross sections.

\subsection{Example: Magnetic fields}

The magnetic moment operator of the neutron is $\muno = \gamma_n \so$, where
$\gamma_n = 2 \mu_n = -3.826$ is the gyromagnetic ratio (spin and magnetic
moment is anti-parallel as for an electron)~\footnote{Note that if we had used
S (with values $S=\pm 1$) to define $\gamma_n$ we would get $\gamma_n =
-1.913$ which is also commonly used.}

A magnetic field, $\BB$, will exert a torque, $\tauB = d\sB/dt = (1/\gamma_n)d\muB/dt$, on the neutron magnetic moment:

\begin{equation}
  \label{eq:torque}
  \frac{1}{\gamma_n} \frac{d\muB}{dt} = \muB \times \BB
\end{equation}

The magnetic moment $\mu$ can be related to the polarization as $\muB
= \gamma_n \PB/2$, and inserting in Eq.~\ref{eq:torque} we find:

\begin{equation}
  \label{eq:pol_magnetic}
  \frac{d\PB}{dt} = \gamma_n \PB \times \BB
\end{equation}

In the simple case where $\BB = (0, 0, B)$, we find the solution
(\cite{gavin} p.~18) :

\begin{eqnarray}
  \nonumber
  P_X(t) & = & \cos(\omega_L t) P_X(0) - \sin(\omega_L t) P_Y(0) \\
  \label{eq:precession}
  P_Y(t) & = & \sin(\omega_L t) P_X(0) + \cos(\omega_L t) P_Y(0) \\
  \nonumber
  P_Z(t) & = & P_Z(0),
\end{eqnarray}

where $\omega_L = -\gamma_n B/\hbar$ is the Larmor frequency.\\

\begin{quote}
  The equations above were checked against the equations in the ``polarimetrie
  neutronique'' notes by Francis Tasset and found to be consistent. There can
  be sign differences between different publications depending on whether they
  use a right-handed (like e.g. McStas) or a left-handed (like e.g. NISP)
  coordinate system.
\end{quote}

This closed-form solution is exactly what the \textbf{Pol\_constBfield} and
\textbf{Pol\_FieldBox} components (section~\ref{sec:new}) evaluate directly
in one step for a spatially uniform field. For the general case of a
spatially varying, rotating, or otherwise non-uniform field, \MCS instead
integrates the precession numerically -- this general algorithm is
described in section~\ref{sec:precession}, and is exactly the same physics
applied piecewise over short enough intervals that the field can be treated
as locally uniform.

\section{Polarized Neutron Scattering}
\label{sec:scat}

First we will give a short introduction to how calculations are done
and then quote some results which are important for implementing the
first McStas components.

All the potentials (nuclear, magnetic, and electric) we will be interested in
can be written on the form:
\begin{equation}
  \label{eq:general_pot}
  \hat{v} = \betao + \alphao \cdot \sigmao
\end{equation}

The first term does not affect the spin, while the second term can
change the spin. Let us just remind here that:

\begin{equation}
  \label{eq:pauli_rules}
  \begin{matrix}
    \sigmaH_x \chiU = \chiD, &
    \sigmaH_y \chiU = i\chiD, &
    \sigmaH_z \chiU = \chiU, \\
    \sigmaH_x \chiD = \chiU, &
    \sigmaH_y \chiD = -i\chiU, &
    \sigmaH_z \chiD = -\chiD.
  \end{matrix}
\end{equation}

So that the interaction proportional to $\sigmaH_x$ and $\sigmaH_y$
results in spin flips, while the interactions with $\sigmaH_z$
conserves the spin.

It turns out to be smart to define a density matrix operator:
\begin{equation}
  \rhoo = \chi \chi^\dagger =
  \left( \begin{matrix}
    |a|^2 & ab^\ast \\
    ba^\ast & |b|^2
  \end{matrix} \right)
  = \frac{1}{2}({\cal I}+\PB \cdot \sigmao),
\end{equation}
where $\chi$ is the neutron wave function (Eq.~\ref{eq:neutron_wave}),
and ${\cal I}$ is the unit matrix.

Using the density matrix the elastic cross section can be written as
(\cite{lovesey84}, Eq.~10.31):
\begin{equation}
  \label{eq:master_sigma}
  \frac{d\sigma}{d\Omega} = \text{Tr} \rhoo \hat{v}^\dagger \hat{v}
  = \sum_{\lambda, \lambda'} p_\lambda \text{Tr} \rhoo
  \langle \lambda | \hat{V}^\dagger(\Q) | \lambda' \rangle
  \langle \lambda' | \hat{V}(\Q) | \lambda \rangle
  \delta(E_\lambda - E_{\lambda'}),
\end{equation}

where $\hat{V}$ is the interaction potential and it is understood that
the trace is to be taken with respect only to the neutron spin
coordinates. The outgoing polarization is given as:
\begin{equation}
  \label{eq:master_pol}
  \PB' \frac{d\sigma}{d\Omega} =
  \text{Tr} \rhoo \hat{v}^\dagger \sigmao \hat{v}
  = \sum_{\lambda, \lambda'} p_\lambda \text{Tr} \rhoo
  \langle \lambda | \hat{V}^\dagger(\Q) | \lambda' \rangle
  \sigmao
  \langle \lambda' | \hat{V}(\Q) | \lambda \rangle
  \delta(E_\lambda - E_{\lambda'})
\end{equation}

Inserting Eq.~\ref{eq:general_pot} in Eq.~\ref{eq:master_sigma} and
Eq.~\ref{eq:master_pol} results in the two master equations for
polarized neutron scattering:

\begin{equation}
  \label{eq:general_sigma}
  \text{Tr} \rhoo \hat{v}^\dagger \hat{v} =
  \alphao^\dagger \cdot \alphao + \betao^\dagger \betao+
  \betao^\dagger (\alphao \cdot \PB) +
  (\alphao^\dagger \cdot \PB) \betao +
  i\PB \cdot (\alphao^\dagger \times \alphao),
\end{equation}

and

\begin{equation}
  \label{eq:general_pol}
  \text{Tr} \rhoo \hat{v}^\dagger \sigmao \hat{v} =
  \betao^\dagger \alphao
  + \alphao^\dagger \betao
  + \betao^\dagger \betao \PB
  + \alphao^\dagger (\alphao \cdot \PB)
  + (\alphao^\dagger \cdot \PB) \alphao
  - \PB (\alphao^\dagger \cdot \alphao)
  - i \alphao^\dagger \times \alphao
  + i \betao^\dagger (\alphao \times \PB)
  + i (\PB \times \alphao^\dagger) \betao.
\end{equation}

Based on these two equations and the interaction potentials all the
results presented in the following are derived in~\cite{lovesey84}.

\subsection{Example: Nuclear scattering}

The nuclear scattering potential for a crystal is:
\begin{equation}
  \label{eq:nuclear_pot}
  \hat{V}_N(\Q) = \sum_{\lB,\dB} \exp(i\Q \cdot \RB_{ld})(A_{ld} +
  \frac{1}{2} B_{ld} \sigmao \cdot \Io_{ld}),
\end{equation}

so that

\begin{eqnarray}
  \label{eq:ab_nuclear}
  \alphao & = &
  \sum_{\lB,\dB} \exp(i\Q \cdot \RB_{ld}) \frac{1}{2} B_{ld} \Io_{ld} \\
  \betao & = &
  \sum_{\lB,\dB} \exp(i\Q \cdot \RB_{ld}) A_{ld},
\end{eqnarray}

where $\Io$ is the nuclear spin operator and the constants $A$ and $B$
are related to the nuclear scattering lengths $b^+$ and $b^-$ as
$A=((I+1)b^++Ib^-)/(2I+1)$ and $B=(b^++b^-)/(2I+1)$.

To calculate the polarization cross section and outgoing polarization
we have to average over the nuclear spin (which we assume is random
oriented), so that terms linear in $\alphao$ (three last terms in
Eq.~\ref{eq:general_sigma}) disappears. The scattering cross section
ends up being (see~\cite{lovesey84} p.~159):

\begin{equation}
  \label{eq:nuclear_sigma}
  \frac{d\sigma}{d\Omega}  =
  \sum_{\lB, \dB, \lB', \dB'} \exp(i \Q \cdot (\RB_{ld}-\RB_{l'd'}))
  (\madsq + \delta_{\lB,\lB'}\delta_{\dB,\dB'}[\sqmad - \madsq
  + \frac{1}{4}\bd])
\end{equation}

where the first term is the coherent cross-section and the second term is the
site-incoherent cross-section. Both terms are independent of $\PB$ as expected
for a system without a preferred internal direction.

The polarization in the final state is:

\begin{equation}
  \label{eq:nuclear_pol}
  \PB' \frac{d\sigma}{d\Omega}  =
  \sum_{\lB, \dB, \lB', \dB'} \exp(i \Q \cdot (\RB_{ld}-\RB_{l'd'}))
  \PB' (\madsq + \delta_{\lB,\lB'}\delta_{\dB,\dB'}[\sqmad - \madsq
  - \frac{1}{12}\bd])
\end{equation}

Comparing Eq.~\ref{eq:nuclear_sigma} and Eq.~\ref{eq:nuclear_pol} we
find that: 1) The nuclear coherent polarization is the same as the
initial polarization. 2) The same is true for the incoherent
scattering due to the random isotope distribution. 3) The nuclear
incoherent scattering due to the random nuclear spin orientations has
polarization $\PB' = -1/3 \PB$ (for a random nuclear spin the
associated Pauli matrix will 2/3 of the time point in the direction of
$\sigmaH_x$ and $\sigmaH_y$ which according to
Eq.~\ref{eq:pauli_rules} flips the spin). \\

For Vanadium, where there is only one isotope and coherent scattering
is negligible, we find $\PB' = -1/3\PB$. There is however one
catch. If the probability for multiple scattering is large one has to
take into account that after two scattering one has: $\PB'(2) =
1/9\PB$, and so forth. The average polarization after a thick vanadium
target is therefore a sum of different contributions.

\subsection{Example: Polarizing Monochromator and Guides}
\label{sub:mono}

\begin{figure}[htbp]
  \begin{center}
    \includegraphics[keepaspectratio,
    width=0.7\columnwidth]{figures/monochromator_pol}
    \caption{Principle and geometry of a polarizing monochromator.}
    \label{fig:mono_princip}
  \end{center}
\end{figure}

In a polarized monochromator and polarizing guides we have a ferromagnetic
crystal in an external magnetic field. The scattering potential is now both
nuclear (no internal direction) and magnetic (internal direction), so in
general the outgoing polarization can be quite complex. However, as
illustrated in Figure~\ref{fig:mono_princip}, the typical setup has many
geometrical constraints: $\tN \cdot \tQ = 0$, $\tN \cdot \PB_\perp = \tN \cdot
\PB$, and $\Q \times (\tN \times \Q) = \tN$, which simplifies the problem.

In~\cite{lovesey84} the calculation for a centrosymmetric ferromagnetic
crystal is done, and inserting the constraints above one finds
(\cite{lovesey84}, Eq.~10.96 and Eq.~10.110):

\begin{eqnarray}
  \label{eq:mono_sigma}
  d\sigma/d\Omega & = & \FN^2 + 2\FN\FM (\PB \cdot \tN) + \FM^2\\
  \label{eq:mono_pol}
  \PB' d\sigma/d\Omega & =
  & \PB [\FN^2 - \FM^2] + \tN [2\FN\FM + 2(\tN \cdot \PB)\FM^2]
\end{eqnarray}

\begin{quote}
  NB! Note that in~\cite{gavin} Eq.~2.2.25 there is a minus in front of the
  second term in Eq.~\ref{eq:mono_sigma}. We have not been able to understand
  this discrepancy, which is probably due to notation. Most other authors
  agree with the minus in front of the second term (e.g. Squires and Francis
  Tasset).
\end{quote}

For a beam which is initially unpolarized we find the outgoing
polarization to be:
\begin{equation}
  \PB' = \frac{\tN 2\FN \FM}{d\sigma/d\Omega}
  = \frac{2\FN\FM}{\FN^2 + \FM^2} \tN,
\end{equation}

so that the beam is fully polarized along $\tN$ if $\FN = \pm \FM$. \\

What we use to characterize the polarizing monochromator in practice
is not $\FN$ and $\FM$, but instead the reflection probabilities $\Ru$
and $\Rd$ (for the reflection of interest).

If we assume that the reflection probabilities are directly proportional to
the cross sections (with proportionality constant $k$), i.e., $\Ru = k
d\sigma/d\Omega(\PB=+\tN)$ and $\Rd = k d\sigma/d\Omega(\PB=-\tN)$ then we can
use Eq.~\ref{eq:mono_sigma} to determine $\FN$ and $\FM$:

\begin{eqnarray}
  \label{eq:mono_up}
  \Ru & = & k (\FN + \FM)^2,\\
  \label{eq:mono_down}
  \Rd & = & k (\FN - \FM)^2.
\end{eqnarray}

The values of $\sqrt{k}\FN$ and $\sqrt{k}\FM$ are then between -1 and +1 and
unit less like the reflection probabilities. In the following we ignore $k$ and
just talk about $\FN$ and $\FM$.

In principle there are four solutions for $\FN$ and $\FM$, so in the code we
currently choose the values where $\FN + \FM = +\sqrt{\Ru}$ and $\FN - \FM =
+\sqrt{\Rd}$ (so that $\FN>0$ and $\FN>\FM$). We then find:

\begin{eqnarray}
  \label{eq:mono_nuc}
  \FN & = & \frac{\sqrt{\Ru} + \sqrt{\Rd}}{2},\\
  \label{eq:mono_mag}
  \FM & = & \frac{\sqrt{\Ru} - \sqrt{\Rd}}{2}.
\end{eqnarray}

When $\FN$ and $\FM$ are determined from these equations,
Eq.~\ref{eq:mono_sigma} and Eq.~\ref{eq:mono_pol} can easily be used to handle
any situation.

This solution is both used for monochromators and guides.

It is not clear that this solution is correct. If we make a simple example
with $\Ru = 1$ and $\Rd = 0.25$ then we could in principle have four
solutions, but let us just quote the two where $\FN$ is positive since the
last two are found by inserting a minus before all the solutions and this does
not change the physics. The two solutions are $\FN = 0.75, \FM=0.25$ and $\FM
= 0.75, \FN=0.25$. All solutions gives the same cross section, but if the
incoming beam is polarized (and only then) the outgoing beam will have two
different polarization values, since $\PB [\FN^2 - \FM^2]$ and $\tN 2(\tN
\cdot \PB)\FM^2$ are different for the two solutions. It seems that one needs
some additional information to choose between the two solutions.

\begin{quote}
  NB! The simplifying geometry shown in Figure~\ref{fig:mono_princip} only
  applies for the sides of the guide wall and not the top and bottom (assuming
  that the magnetizing field is pointing up or down), so there another set of
  equations should really be used.
\end{quote}

The same physics could also be used for a polarizing powder or single crystal
sample if $\FN$ and $\FM$ can be calculated with some other program, but one
would have to use the general form of Eq.~\ref{eq:mono_sigma} and
Eq.~\ref{eq:mono_pol} without the simplifying geometrical constraints for
monochromators and guides.


%%%%%%%%%%%%%%%%%%%%%%%%%%%%%%%%%%%%%%%%%%%%%%%%%%%%%%%%%%%%%%%%%%%%%%%%%%%%%%%%
\section{Magnetic fields: the precession algorithm and field framework}
\label{sec:precession}
\index{Polarisation!precession algorithm|textbf}
\index{Library!pol-lib}

All magnetic-field-related polarisation components in \MCS share a common
run-time library, \texttt{pol-lib} (\texttt{mccode/nlib/share/pol-lib.c}\slash
\texttt{.h}), which provides the spin-precession integrator and a small set
of built-in field functions. A separate library, \texttt{ref-lib}, provides
the standard supermirror reflectivity function (\texttt{StdReflecFunc},
section~\ref{sub:mono}) used by the polarising mirrors and guides.

\subsection{The magnetic field stack}
\label{sec:magnetstack}

Rather than each field-producing component computing precession locally and
independently, \texttt{pol-lib} maintains a small per-particle \emph{stack}
of active magnetic fields, \texttt{\_particle->mcMagnet}. A component that
defines a field region (e.g.\ \textbf{Pol\_Bfield}) \emph{pushes} a
description of its field onto this stack when a particle enters the region,
using
\begin{quote}
\texttt{mcmagnet\_push(\_particle, field\_type, \&rotation, \&position, stopbit, params[8])}
\end{quote}
where \texttt{field\_type} selects one of the built-in field functions
(section~\ref{sec:fieldfuncs}), \texttt{rotation}/\texttt{position} record
the field region's placement in the lab frame (so the field can later be
evaluated in the region's own local coordinates regardless of where in the
instrument it sits), and \texttt{params} holds up to 8 doubles of
field-specific parameters (amplitude, extent, \ldots). A companion
\textbf{Pol\_Bfield\_stop} component calls \texttt{mcmagnet\_pop()} to
remove the field from the stack once the particle leaves the region.

Because this is a stack rather than a single active field, \textbf{field
regions may be nested or overlap}: any component placed between a
\texttt{Pol\_Bfield}/\texttt{Pol\_Bfield\_stop} pair -- including another,
independent field region -- experiences the vector sum of every field
currently on the stack (evaluated by \texttt{mcmagnet\_get\_field()}, which
walks the whole stack, transforms each field into the lab frame, and adds
the contributions). A field region may also be declared \emph{closed}
(\texttt{concentric=0}), in which case it pushes and pops its own field
internally and no other components may be placed inside it -- this is the
simpler, non-nestable mode used when a single self-contained field region
is all that is needed.

Whenever a particle is propagated through a region with an active field
(technically: whenever \texttt{PROP\_DT} is called while
\texttt{\_particle->mcMagnet} is non-\texttt{NULL}), the propagation macro
transparently invokes the precession integrator described next, so that
ordinary component code (slits, monitors, samples placed inside a field
region) does not need to know anything about polarisation -- the spin
precesses correctly regardless of what else is happening to the particle
during that propagation step.

\subsection{Numerical spin precession: \texttt{SimpleNumMagnetPrecession}}
\label{sec:numprecession}
\index{Polarisation!SimpleNumMagnetPrecession}

The general-purpose precession routine, \texttt{SimpleNumMagnetPrecession()},
integrates Eq.~\ref{eq:pol_magnetic} numerically along the particle's
straight-line trajectory (gravity is \emph{not} accounted for inside the
field itself) for a total propagation time $dt$, using an
\emph{adaptive-step} scheme:

\begin{enumerate}
\item The field $\mathbf{B}_\mathrm{start}$ is evaluated at the particle's
  current position (by querying the whole field stack, or -- for the
  self-contained field components -- a single field function).
\item A trial sub-step of length \texttt{mc\_pol\_initial\_timestep}
  (default $10^{-5}$~s, capped at the remaining $dt$) is proposed, and the
  field $\mathbf{B}_\mathrm{end}$ is evaluated at the particle's position at
  the end of that trial step.
\item If the angle between $\mathbf{B}_\mathrm{start}$ and
  $\mathbf{B}_\mathrm{end}$ exceeds a threshold, \texttt{mc\_pol\_angular\_accuracy}
  (default $1^\circ$), the trial step is halved and re-evaluated; this
  repeats until the field direction changes by less than the threshold
  over the sub-step (or the step size underflows to \texttt{FLT\_EPSILON}).
  This is what makes the scheme \emph{adaptive}: in a slowly-varying or
  uniform field, the integrator quickly settles on large sub-steps (up to
  the $10^{-5}$~s cap), while in a rapidly rotating field (e.g. the
  \texttt{rotating} or \texttt{majorana} functions of
  section~\ref{sec:fieldfuncs}) it automatically refines the step until the
  field is well approximated as locally uniform.
\item Once an acceptable sub-step is found, the particle position and time
  are advanced by that sub-step, and the spin is precessed \emph{about the
  mean field over the sub-step}, $\bar{\mathbf{B}} = \tfrac12
  (\mathbf{B}_\mathrm{start}+\mathbf{B}_\mathrm{end})$, by the angle
  \[
    \phi = \left(|\bar{\mathbf{B}}|\, \delta t\, \omega_L\right) \bmod 2\pi,
    \qquad
    \omega_L = -2\pi \times 29.16~\mathrm{MHz/T},
  \]
  matching Eq.~\ref{eq:precession} applied over the sub-step, using a
  generic axis-angle rotation of $(s_x,s_y,s_z)$ about $\bar{\mathbf{B}}$
  (rather than the special-case $z$-axis form of Eq.~\ref{eq:precession},
  since $\bar{\mathbf{B}}$ may point in any direction).
\item This is repeated, consuming the remaining $dt$ sub-step by sub-step,
  until the whole requested propagation time has been precessed over.
\end{enumerate}
Positions/velocities/spin are tracked internally in the lab frame (fields
are pushed to the stack together with the rotation/position needed to
transform into and out of each field region's local frame) and the final
spin is transformed back into the calling component's local frame before
being returned. \textbf{Pol\_tabled\_field} (section~\ref{sec:new})
re-implements this identical algorithm internally (rather than going
through the shared field stack), so that it can query its own tabulated/
interpolated field directly.

Two component-specific accuracy parameters, exposed via
\texttt{mc\_pol\_set\_timestep()} and \texttt{mc\_pol\_set\_angular\_accuracy()},
allow the initial sub-step size and the angular-accuracy threshold to be
tuned instrument-wide if the defaults are not adequate for a particular
field geometry.

\subsection{Available field functions}
\label{sec:fieldfuncs}
\index{Polarisation!field functions}

The field pushed onto the stack by \textbf{Pol\_Bfield} is one of the
following built-in functions, selected by the integer \texttt{field\_type}
parameter (\texttt{Pol\_FieldBox} and \texttt{Pol\_constBfield} always use
the constant-field case, evaluated in closed form rather than via the
stack):

\begin{description}
\item[1 -- constant] $\mathbf{B}=(B_x,B_y,B_z)$, uniform throughout the
  region.
\item[2 -- rotating] $\mathbf{B}$ starts as $(0,B_y,0)$ at the entrance and
  rotates smoothly to $(B_y,0,0)$ at the far end (over the region's
  \texttt{zdepth}), following $B_x = B\sin(\tfrac{\pi}{2}z/L)$, $B_y =
  B\cos(\tfrac{\pi}{2}z/L)$.
\item[3 -- majorana] a large component along $x$ decreases linearly from
  $+B_x$ to $-B_x$ across the region (passing through zero at the centre --
  the classic ``Majorana flip'' geometry for studying spin-flip loss at a
  field zero-crossing), plus a small, constant transverse component along
  $y$.
\item[6 -- gradient] linear interpolation of the full field vector between
  two given endpoint vectors $\mathbf{B}_1$ (at $z=-L/2$) and
  $\mathbf{B}_2$ (at $z=+L/2$) -- a generalisation of the majorana case to
  an arbitrary field-vector gradient.
\item[-1 -- tabled] used internally by \textbf{Pol\_tabled\_field}: the
  field is not one of the analytic functions above but is instead
  interpolated (via \texttt{interpolation-lib}, choosing between a
  \texttt{kdtree} or \texttt{regular}-grid interpolator) from an externally
  supplied point cloud file of $(x,y,z,B_x,B_y,B_z)$ rows -- this is the
  route to use for a realistic, measured or finite-element-computed field
  map.
\item[4 (MSF) and 5 (RF)] reserved identifiers for a modulated/RF spin
  flipper field; not implemented in the current release (a radio-frequency
  field is more commonly modelled explicitly via a rotating-frame field
  function or, for an idealised flipper, by simply mirroring the spin
  vector, see \textbf{Pol\_SF\_ideal} below).
\end{description}
New field functions can be added by writing a function with the signature
of e.g.\ \texttt{const\_magnetic\_field()} in \texttt{pol-lib.c} and adding
a case to \texttt{magnetic\_field\_dispatcher()}; per the file header, this
requires modifying the shared library rather than being an
end-user-facing extension point.

%%%%%%%%%%%%%%%%%%%%%%%%%%%%%%%%%%%%%%%%%%%%%%%%%%%%%%%%%%%%%%%%%%%%%%%%%%%%%%%%
\section{Polarisation components in \MCS}
\label{sec:new}

The components below can be grouped by role:
\begin{itemize}
\item \textbf{Setting and analysing polarisation} (unphysical, for testing/idealised use).
\item \textbf{Polarising monochromators, mirrors and guides} (physical, reflectivity-based).
\item \textbf{Magnetic field regions} (precession, using the framework of section~\ref{sec:precession}).
\item \textbf{Polarisation-sensitive monitors}.
\item \textbf{Samples} affecting polarisation.
\end{itemize}
As of \MCS~3.x these are found in the \texttt{optics} and \texttt{monitors}
component categories (not a dedicated ``polarisation'' category); a
component is polarisation-aware precisely when its \texttt{TRACE} section
reads and/or writes the particle's spin state $(s_x,s_y,s_z)$.

\subsection{Setting and analysing polarisation}

\begin{itemize}
\item \textbf{Set\_pol:} An unphysical, zero-size component (drawn like an
  \texttt{Arm}) used to impose a polarisation state, in one of four modes
  selected by \texttt{randomOn} and \texttt{normalize}: (i) hard-code the
  polarisation to $(p_x,p_y,p_z)$; (ii) set it to a random vector \emph{on}
  the unit sphere (fully polarised, random direction); (iii) set it to a
  random vector \emph{within} the unit sphere (random direction and
  degree of polarisation); (iv) hard-code the \emph{direction} of
  $(p_x,p_y,p_z)$ but normalise it to unit polarisation.

\item \textbf{PolAnalyser\_ideal:} An unphysical, ideal polarisation
  analyser. Given an analysis direction $\mathbf{m}=(m_x,m_y,m_z)$
  ($|\mathbf{m}|\le 1$ for an imperfect analyser), it reduces the neutron
  weight by the projection probability
  $\tfrac12\bigl(1+\mathbf{S}\cdot\mathbf{m}\bigr)$ and collapses the
  outgoing spin to $\mathbf{m}$, or absorbs the ray if that probability is
  zero or negative.

\item \textbf{Pol\_SF\_ideal:} An idealised spin flipper: within an
  absorbing box, the polarisation vector is mirrored through the plane
  through the origin with normal $(n_x,n_y,n_z)$ ($\mathbf{S}' = \mathbf{S}
  - 2(\mathbf{S}\cdot\hat{\mathbf{n}})\hat{\mathbf{n}}$); rays missing the
  box pass through untouched. This is the natural building block for a
  spin-echo instrument's $\pi$-flip coils (section~\ref{sec:spinecho}) when
  an idealised, 100\%-efficient flipper is sufficient.
\end{itemize}

\subsection{Polarising monochromators, mirrors and guides}
\label{sub:mono-comps}

These all implement the physics of section~\ref{sub:mono}
(Eqs.~\ref{eq:mono_sigma}--\ref{eq:mono_mag}): given reflectivities for
spin up and down ($R_\uparrow$, $R_\downarrow$, either as
\texttt{\{R0,Qc,alpha,m,W\}} parametrised curves via \texttt{ref-lib}'s
\texttt{StdReflecFunc}, or as tabulated reflectivity files), the shared
helper functions \texttt{GetMonoPolFNFM}, \texttt{GetMonoPolRefProb},
\texttt{SetMonoPolRefOut} and \texttt{SetMonoPolTransOut} (all in
\texttt{pol-lib}) compute $F_N$, $F_M$, the reflection probability, and the
outgoing polarisation for both the reflected and transmitted branches.
Since a 2024 revision (E.B.~Knudsen, P.~Willendrup, H.~Lee), all of these
components treat a polarising reflection/transmission event as a
projective quantum measurement along the mirror's quantisation axis: the
in-plane spin components are explicitly zeroed ($s_x=s_z=0$) at the point
of interaction, before the new $s_y$ is set by \texttt{SetMonoPolRefOut}/
\texttt{SetMonoPolTransOut} -- reflecting that only the component of
polarisation along the mirror's up direction survives a measurement-like
interaction of this kind.

\begin{itemize}
\item \textbf{Monochromator\_pol:} A flat, infinitely thin mosaic
  monochromator/analyser crystal (Gaussian mosaic and $d$-spread,
  billiard-ball reflection), the direct polarising analogue of
  \textbf{Monochromator\_flat}, parametrised by \texttt{Rup}/\texttt{Rdown}
  (or \texttt{Q}/\texttt{DM} for the lattice spacing).
\item \textbf{Pol\_mirror:} A flat, infinitely thin, non-absorbing
  polarising mirror in the $y$-$z$ plane, reflecting and/or transmitting
  according to \texttt{p\_reflect} (use $-1$ to sample physically from the
  mirror's own reflectivity, or a fixed value to force a chosen
  reflect/transmit statistics split while re-weighting correctly).
\item \textbf{Pol\_bender:} A curved, multi-slit polarising bender (based
  on \textbf{Guide\_curved}), with independently specifiable
  reflectivities for each of the four walls (top/bottom/left/right) and
  gravity support.
\item \textbf{Pol\_guide\_mirror}\ /\ \textbf{Pol\_guide\_vmirror:}
  Straight rectangular guides with one (\texttt{\_mirror}) or two,
  V-shaped (\texttt{\_vmirror}) polarising supermirrors sitting on the
  diagonal(s) inside an otherwise ordinary (optionally non-polarising)
  guide. Setting the up and down reflectivities of the diagonal mirror
  equal turns either component into an unpolarised \emph{frame-overlap
  mirror} instead -- a deliberate dual use of the same geometry.
\end{itemize}

\subsection{Magnetic field components}
\label{sub:fieldcomps}

See section~\ref{sec:precession} for the shared precession algorithm and
field functions used by these.

\begin{itemize}
\item \textbf{Pol\_Bfield} / \textbf{Pol\_Bfield\_stop:} The general-purpose
  field region, supporting all of the field functions of
  section~\ref{sec:fieldfuncs} and a box, cylindrical, spherical, or
  window-shaped extent. Used together as a \texttt{concentric} pair
  (\texttt{Pol\_Bfield} pushes the field, other components may be placed
  freely inside, \texttt{Pol\_Bfield\_stop} pops it), or standalone as a
  closed region (\texttt{concentric=0}) that nothing else may be placed
  inside.
\item \textbf{Pol\_FieldBox:} A simple, self-contained, non-nestable box
  with a single uniform field $(B_x,B_y,B_z)$, evaluated with a direct,
  single-step closed-form precession (Eq.~\ref{eq:precession}) rather than
  the general stack/integrator -- the lightest-weight choice for a plain
  guide field or spin-flip coil.
\item \textbf{Pol\_tabled\_field:} A field region (box/cylinder/sphere, or
  an arbitrary OFF/PLY-geometry volume) whose field is interpolated from an
  externally supplied point-cloud file, using the same adaptive-step
  precession algorithm re-implemented internally
  (section~\ref{sec:numprecession}). The interpolation backend
  (\texttt{kdtree} or \texttt{regular}-grid) can be chosen explicitly or
  left to an automatic CPU/GPU-appropriate default.
\item \textbf{Pol\_constBfield:} A rectangular, closed (non-nestable) box
  with a constant field along its local $y$-axis, evaluated with the
  closed-form solution of Eq.~\ref{eq:precession} directly (no field
  stack). Convenient for guide fields and simple spin flippers: rather
  than specifying $B$ directly, \texttt{fliplambda}/\texttt{flipangle} let
  the field be specified as ``the field that flips a neutron of this
  wavelength by this angle over this component's length'' -- exactly the
  \texttt{GetConstantField()} helper of \texttt{pol-lib}.
\end{itemize}

\subsection{Polarisation-sensitive monitors}

\begin{itemize}
\item \textbf{Pol\_monitor:} Measures the projection of the polarisation
  onto a user-defined direction $\mathbf{m}=(m_x,m_y,m_z)$, i.e.\
  $\mathbf{m}\cdot\mathbf{S}$, integrated over all detected rays.
\item \textbf{PolLambda\_monitor:} As \textbf{Pol\_monitor}, resolved as a
  function of wavelength $\lambda$.
\item \textbf{MeanPolLambda\_monitor:} As \textbf{PolLambda\_monitor}, but
  reporting the \emph{mean} polarisation projection per wavelength bin
  (rather than the summed/integrated signal) -- the natural monitor for a
  spin-echo polarisation-vs.-parameter scan (section~\ref{sec:spinecho}).
\item \textbf{PSD\_monitor\_4PI\_spin:} A spherical, $4\pi$
  position-sensitive monitor that additionally records the mean
  polarisation projection per pixel, for mapping how polarisation varies
  with scattering angle.
\end{itemize}

\subsection{Samples and polarisation}

Few \MCS sample components currently propagate polarisation explicitly.
The generic \textbf{Incoherent} component (the modern, general-purpose
replacement for the historical, now-obsolete \textbf{V\_sample}) can be
parametrised to reproduce the classic depolarising-incoherent-scatterer
result of section~\ref{sec:scat}'s nuclear-scattering example
($\mathbf{P}' = -\tfrac13\mathbf{P}$ per single incoherent scattering
event, from the random projection of nuclear spin onto $\hat\sigma_x$ and
$\hat\sigma_y$); as noted there, this simple factor is only exact for a
single scattering event, and multiple scattering must be handled with
some care (consecutive depolarisation by $(-\tfrac13)^n$ is only valid if
each scattering order can be isolated). Magnetic (Bragg) scattering from
an ordered magnetic structure, using the magnetic structure factor
$F_M$ of section~\ref{sec:pol}, is not yet implemented as a general sample
component; \textbf{Single\_crystal} and \textbf{PowderN} presently model
only nuclear scattering and do not read or write the spin state.

%%%%%%%%%%%%%%%%%%%%%%%%%%%%%%%%%%%%%%%%%%%%%%%%%%%%%%%%%%%%%%%%%%%%%%%%%%%%%%%%
\section{Test and example instruments}
\label{sec:test}

Every instrument below is found in the \texttt{Tests\_polarization}
category of the example library (browsable via \texttt{mcgui}'s
``New from Template'' menu, or \texttt{mcdoc}).

\subsection{Component-level regression tests}
These verify a polarising component reduces correctly to its
non-polarising counterpart, or exercise one component's specific physics
in isolation: \textit{Test\_Pol\_Bender}, \textit{Test\_Pol\_FieldBox},
\textit{Test\_Pol\_Guide\_mirror}, \textit{Test\_Pol\_Guide\_Vmirror},
\textit{Test\_Pol\_Mirror}, \textit{Test\_Pol\_MSF},
\textit{Test\_Pol\_SF\_ideal}, \textit{Test\_Pol\_Set},
\textit{Test\_Pol\_Tabled}, \textit{Test\_pol\_ideal},
\textit{Test\_Magnetic\_Constant}, \textit{Test\_Magnetic\_Majorana}, and
\textit{Test\_Magnetic\_Rotation} (illustrating, respectively, the
constant, majorana, and rotating field functions of
section~\ref{sec:fieldfuncs}).

\subsection{Test\_Pol\_TripleAxis}
A complete example of a polarised triple-axis spectrometer: polarising
monochromator and analyser, a vanadium (depolarising incoherent) sample,
and a spin flipper -- illustrating how the individual components of
section~\ref{sec:new} combine into a real instrument.

\subsection{He3\_spin\_filter}
Illustrates the use of a $^3$He spin filter cell for beam or analysis-side
polarisation.

\subsection{A spin-echo example: SE\_example}
\label{sec:spinecho}
\index{Spin-echo}

\textit{SE\_example} (and its variant \textit{SE\_example2}, both written
for the 2010 PNCMI school in Delft) is a mock-up transmission spin-echo
instrument, and a compact illustration of most of the machinery described
in this chapter working together:
\begin{itemize}
\item \textbf{Pol\_mirror} polarises the incoming beam and, symmetrically,
  analyses the scattered beam.
\item \textbf{Pol\_Bfield} regions define the two (oppositely-signed)
  precession guide fields on either side of the sample position, exactly
  the nested/concentric field-stack usage of section~\ref{sec:magnetstack}
  (other components -- the sample, monitors -- sit \emph{inside} each
  field region, between its \texttt{Pol\_Bfield}/\texttt{Pol\_Bfield\_stop}
  pair).
\item \textbf{MeanPolLambda\_monitor} records the mean polarisation
  projection as a function of wavelength -- the classic spin-echo readout,
  whose oscillation as a small guide-field perturbation \texttt{dBz} is
  scanned traces out the echo signal.
\end{itemize}
A typical scan takes the form
\begin{lstlisting}
mcrun SE_example.instr dBz=-0.0001,0.0001 -N41 -n1e5
\end{lstlisting}
sweeping a small field imbalance \texttt{dBz} between the two precession
regions and recording the resulting polarisation echo via
\textbf{MeanPolLambda\_monitor}. The instrument's own parameters
(\texttt{POL\_ANGLE}, \texttt{Bguide}, \texttt{Bflip}, \texttt{dBz},
\texttt{Lam}, \texttt{dLam}) let the polariser/analyser angle, the guide-
and flipper-field strengths, and the source wavelength band be adjusted
independently, and an optional incoherent scatterer (\texttt{SAMPLE=1})
can be inserted to study the effect of depolarising sample scattering on
the echo.
